% =============================================================================
% Kapitel 4: Konzeption der Informationsverarbeitungsarchitektur
% ARBEITSGERUEST / REVIEW CANDIDATE
% Stand: 03.09.2026
%
% Dieses Dokument enthaelt bewusst NOCH KEINEN ausformulierten Thesis-Text.
% Es bildet die Argumentationsstruktur fuer das spaetere Kapitel 4 ab.
%
% Zielstruktur:
% Kapitel 3 = WIE das Konzept entwickelt und untersucht wurde
% Kapitel 4 = WAS als Architekturkonzept resultiert
% Kapitel 5 = WIE das Konzept prototypisch realisiert wurde
% Kapitel 6 = WIE es verifiziert und validiert wurde
% Kapitel 7 = WAS beobachtet wurde
% Kapitel 8 = WAS die Ergebnisse bedeuten
%
% Dieses Kapitel konsolidiert spaeter die bisherigen Platzhalter:
% - Kapitel/04_StakeholderAnalyse.tex
% - Kapitel/05_Architekturkonzeption.tex
%
% Quellenstatus:
% [REPO-BIB]  Quelle ist im Thesis-Repository und in literatur.bib vorhanden.
% [REPO-FILE] Quelle ist als Datei vorhanden, aber bibliographisch oder
%             claimweise noch zu pruefen.
% [ADD]       Quelle muss vor Verwendung beschafft, archiviert und geprueft werden.
% [MODEL]     Aussage stammt aus dem akzeptierten CATIA SysML-v2-Modell.
% [PROJECT]   Aussage stammt aus akzeptierten ADRs, Checkpoints oder Contracts.
%
% Schreibregel:
% Projektspezifische Architekturentscheidungen duerfen durch MODEL/PROJECT
% getragen werden. Allgemeine wissenschaftliche Aussagen benoetigen Literatur.
% =============================================================================

\section{Konzeption der Informationsverarbeitungsarchitektur}
\label{sec:architekturkonzept}

% -----------------------------------------------------------------------------
% FUNKTION DES KAPITELS
% -----------------------------------------------------------------------------
%
% Zweck:
% Dieses Kapitel beschreibt das Ergebnis der in Kapitel 3 dargestellten
% Konzeptionsmethodik.
%
% Der Leser soll nach diesem Kapitel verstanden haben:
%
% 1. Welches System betrachtet wird und welche Akteure ausserhalb liegen.
% 2. Welche Anforderungen die Architektur treiben.
% 3. Welche wesentlichen Funktionen das System erfuellen muss.
% 4. Warum die finale Architektur in neun fachliche Subsysteme zerlegt ist.
% 5. Wie Engineering Information durch das System fliesst.
% 6. Wo LLMs eingesetzt werden und wo ihre Authority endet.
% 7. Wo Human Engineering Authority und Model Authority liegen.
% 8. Wie Provenance und Traceability ueber die Verarbeitung erhalten werden.
% 9. Wie mehrere Quellen ohne versteckte Source-Prioritaet zusammengefuehrt
%    bzw. gegeneinander bewertet werden.
% 10. Warum Model Assembly, SysML-v2-Generierung, Validierung und Publikation
%     getrennte Verantwortlichkeiten sind.
%
% ZENTRALE ARCHITEKTURERZAEHLUNG:
%
% Project / Source Context
% -> source-local Engineering Information Processing
% -> source-local Human Engineering Review
% -> Approved Engineering Information
% -> exact Project Fit / Project Admission
% -> one Project-level Model Candidate Set
% -> Human Model Candidate / Placement authority
% -> deterministic Internal Engineering Model assembly
% -> deterministic SysML v2 generation
% -> separate generated-artifact validation
% -> Final Human Model Review
% -> publication gate / immutable publication
%
% WICHTIGE ABGRENZUNG:
% - Keine Python-Modul-Dokumentation.
% - Keine vollstaendige Wiedergabe aller 153 Subsystem Requirements.
% - Keine vollstaendige Wiedergabe aller 44 Subsystem Functions.
% - Keine vollstaendige Wiedergabe aller 31 Logical Components.
% - Keine Auflistung aller 35 Cross-SBS Interfaces und 56 Item Flows.
% - Vollstaendige Detailstruktur verbleibt im eingereichten CATIA-Modell.
%
% Die Thesis muss dennoch ohne Oeffnen des .mdzip verstaendlich bleiben.


% =============================================================================
\subsection{Systemkontext und architekturtreibende Anforderungen}
\label{sec:architektur_systemkontext}
% =============================================================================
Die Konzeption der Informationsverarbeitungsarchitektur beginnt mit der
Abgrenzung des Turing Generators als System of Interest. Das System verarbeitet
heterogene Engineering-Bestandsinformationen mit dem Ziel, daraus
nachvollziehbare Architekturinformationen und schließlich formale
SysML-v2-Artefakte abzuleiten.

Die Engineering-Quellen selbst sowie die beteiligten menschlichen Rollen liegen
außerhalb der Systemgrenze. Innerhalb der Systemgrenze liegen dagegen die
Verarbeitung, Bewertung, Aufbereitung und Weitergabe der Engineering-Information
sowie die technischen Mechanismen zur Unterstützung menschlicher
Entscheidungen.

Diese Abgrenzung bildet den Ausgangspunkt für die weitere funktionale und
logische Zerlegung der Architektur. Der Systemkontext beschreibt zunächst
ausschließlich relevante externe Akteure, Informationsquellen und deren
Interaktionen mit dem System. Interne Systemfunktionen, logische Komponenten und
Subsystemgrenzen werden auf dieser Betrachtungsebene noch nicht festgelegt.
Damit werden unterschiedliche Architektursichten entsprechend ihrem jeweiligen
Betrachtungszweck getrennt \cite{ISO42010_2022}.


\subsubsection{System of Interest und externe Akteure}
\label{sec:system_of_interest}

% KERNAUSSAGEN:
%
% 1. System of Interest ist der Turing Generator.
%
% 2. Der Turing Generator verarbeitet heterogene Engineering-
%    Bestandsinformationen zu traceablen SysML-v2-Architekturartefakten
%    unter expliziter menschlicher Review- und Freigabeautoritaet.
%
% 3. Engineering Sources liegen ausserhalb der Systemgrenze.
%
% 4. Vier menschliche Rollen bilden den fachlichen Stakeholder-Kontext:
%
%    ACT_01 System Modeler
%    - kuratiert, strukturiert und integriert Systemmodelle
%    - prueft modellierungsbezogene Ergebnisse
%
%    ACT_02 Domain Expert
%    - klaert domaenenspezifische Mehrdeutigkeiten
%    - validiert Annahmen und fachliche Interpretation
%
%    ACT_03 System Architect
%    - verantwortet Architekturconstraints und Strukturqualitaet
%    - prueft und autorisiert Architektur- bzw. Modellinhalte
%
%    ACT_04 Information Manager
%    - verantwortet Informationsqualitaet, Source Reliability,
%      semantische Konsistenz und Informationsgovernance
%
% [MODEL]
% collaboration/CATIAMSOSA_TextualNotation.txt
% - Stakeholder Context CTX_001
% - Engineering Source
% - Turing Generator System
% - ACT_01 bis ACT_04
%
% [REPO-FILE]
% ISO/IEC/IEEE 42010:2022
% Literatur/ISO-IEC-IEEE 42010_2022-11-00_EN_3400447.pdf
% Relevanz:
% - Entity of Interest
% - Environment
% - Stakeholder
% - Concerns
% - Architecture Description
% TODO: korrekten BibTeX-Eintrag fuer Version 2022 anlegen bzw. pruefen.
%
% [REPO-BIB]
% \cite{SysMLv2}
% Relevanz:
% - SysML v2 als Sprache fuer die modellierte Architektur.
%
% ABBILDUNGSKANDIDAT 4.1:
% CATIA Stakeholder Context View.
% Turing Generator in der Mitte, Engineering Sources sowie vier Human Roles
% ausserhalb.
%
% Diese Abbildung gehoert in den Haupttext, weil sie die Systemgrenze definiert.
Das Architekturmodell unterscheidet neben den zu verarbeitenden
Engineering-Quellen vier menschliche Stakeholderrollen. Sie repräsentieren
unterschiedliche fachliche Verantwortlichkeiten innerhalb des betrachteten
Engineering-Prozesses.

Der \emph{System Modeler} verantwortet insbesondere die Strukturierung,
Überprüfung und Integration von Systemmodellen. Im Kontext des Turing Generators
betrifft dies sowohl die Bewertung von Zwischenergebnissen als auch die
abschließende Betrachtung des resultierenden Modells.

Der \emph{Domain Expert} bringt produkt- beziehungsweise domänenspezifisches
Wissen in den Verarbeitungsprozess ein. Seine fachliche Verantwortung liegt
insbesondere dort, wo Ausgangsinformationen mehrdeutig, unvollständig oder
widersprüchlich sind und eine rein automatisierte Interpretation nicht
ausreichend ist.

Der \emph{System Architect} betrachtet die strukturelle Qualität der
resultierenden Architektur. Zu seinen Verantwortlichkeiten gehören
Architekturconstraints, Modularität, Konsistenz sowie die Einordnung erzeugter
Modellinhalte in einen übergeordneten Systemkontext.

Der \emph{Information Manager} verantwortet die Governance der verarbeiteten
Engineering-Information. Dazu gehören insbesondere Informationsqualität,
Source Reliability, semantische Konsistenz sowie die Anwendung
ontologiebezogener Regeln.

Die vier Rollen bilden damit unterschiedliche fachliche Perspektiven auf die
Verarbeitung derselben Engineering-Information ab. Ihre Trennung ist
insbesondere für die später beschriebenen Review- und Authority-Grenzen
relevant, da nicht jede fachliche Entscheidung durch dieselbe Rolle getroffen
werden muss.

% ABBILDUNG:
% CATIA View IV_TuringGeneratorStakeholderContext
%
% Die Abbildung soll enthalten:
% - Turing Generator als System of Interest
% - Engineering Sources außerhalb der Systemgrenze
% - System Modeler
% - Domain Expert
% - System Architect
% - Information Manager
% - Interaktionen über die Systemgrenze
%
% \begin{figure}[htbp]
%     \centering
%     \includegraphics[width=\textwidth]{Abbildungen/...}
%     \caption{Systemkontext des Turing Generators mit externen
%     Engineering-Quellen und menschlichen Stakeholderrollen}
%     \label{fig:turing_systemkontext}
% \end{figure}

Die Interaktionen über die Systemgrenze besitzen dabei unterschiedliche
Funktionen. Engineering-Quellen stellen die zu verarbeitende
Informationsgrundlage bereit. Die menschlichen Akteure interagieren dagegen mit
dem System, um Modellierungsentscheidungen zu treffen, fachliche Unklarheiten zu
klären, Architekturinhalte zu bewerten oder die Verwendung von
Engineering-Information zu steuern.

Der Systemkontext macht damit bereits eine grundlegende Eigenschaft der
Informationsverarbeitungsarchitektur sichtbar: Die Transformation von
Bestandsinformationen in ein Architekturmodell ist kein vollständig autonomer
Prozess. Automatisierte Verarbeitung und menschliche fachliche Verantwortung
werden vielmehr innerhalb eines gemeinsamen Informationsflusses miteinander
verbunden.


\subsubsection{Architekturtreibende User Needs und Anforderungen}
\label{sec:architekturtreibende_bedarfe}

% Zweck:
% Nicht jede User Need und jedes Requirement wiedergeben.
% Stattdessen Requirement-Cluster zeigen, aus denen die Architekturstruktur
% nachvollziehbar wird.
%
% KERNAUSSAGEN:
%
% A. Verarbeitung heterogener Engineering Information
%    - Data Suitability
%    - Completeness
%    - Consistency
%    - Reliability
%    - Relevance
%    - controlled completion
%
% B. Mehrere Projekte und mehrere Sources
%    - Multiple Project Support
%    - Multiple Source Support
%    - Project Information Isolation
%    - Processing State Transparency
%
% C. Semantische Konsistenz
%    - Technical Vocabulary
%    - Semantic Conflict Detection
%    - Ontology / semantic governance
%
% D. Human Review und Authority
%    - Review Need Transparency
%    - Human Review Decision Association
%    - Approved Information Boundary
%    - fachliche Modellvalidierung
%
% E. Architektur- und Modellqualitaet
%    - Structural Pattern Conformance
%    - Model Reusability
%    - Modular Model Growth
%    - Non Destructive Submodel Integration
%    - Interface Definition Consistency
%    - Larger Context Compatibility
%
% F. Architecture as Code und SysML v2
%    - Architecture Information Derivation
%    - Architecture as Code Provision
%    - SysML v2 Target Representation
%
% G. Traceability und Erklaerbarkeit
%    - Source Traceability
%    - Explanation of Generated Engineering Results
%    - Processing History
%
% [MODEL]
% collaboration/CATIAMSOSA_TextualNotation.txt
%
% BESONDERS GUT ALS REPRASENTATIVE REQUIREMENTS:
% SR_001 Input Data Suitability Assessment
% SR_004 Source Reliability Classification
% SR_005 Source Relevance Classification
% SR_006 Architecture Information Derivation
% SR_019 Source Traceability
% SR_027 Architecture as Code Provision
% SR_028 SysML v2 Target Representation
% SR_030 Project Information Isolation
% SR_038 Multiple Source Support
%
% Fuer Assurance Requirements genaue IDs/Namen vor Verwendung aus dem
% finalen CATIA Snapshot ziehen.
%
% [ADD]
% ISO/IEC/IEEE 29148:2018
% Relevanz:
% - Requirements Engineering
% - Herkunft und Traceability von Anforderungen
%
% TABELLEKANDIDAT 4.1:
% Architekturtreibende Requirement-Cluster
%
% Spalten:
% Bedarfskategorie | Representative Requirements | Architekturwirkung
%
% Vollstaendige Requirement-Bevoelkerung NICHT in den Haupttext.


Die fachliche Grundlage der Architektur bilden die im Modell erfassten
User Needs. Sie beschreiben zunächst den aus Sicht der Stakeholder bestehenden
Bedarf, ohne bereits eine konkrete technische Lösung vorzugeben. Aus ihnen
wurden Stakeholder Requirements abgeleitet, die den Bedarf in Anforderungen an
das System überführen.

Die vollständige Anforderungsstruktur wird im CATIA-Modell über mehrere
Abstraktionsebenen weitergeführt. Eine vollständige Wiedergabe aller User Needs
und Requirements im Haupttext würde jedoch weitgehend den Inhalt des
Architekturmodells duplizieren. Tabelle~\ref{tab:architekturtreibende_requirements}
fasst deshalb diejenigen Bedarfskategorien zusammen, die für die Gestaltung der
Informationsverarbeitungsarchitektur besonders prägend sind.

\begin{table}[htbp]
\caption{Architekturtreibende User Needs und repräsentative Stakeholder Requirements}
\label{tab:architekturtreibende_requirements}
\centering
\small
\begin{tabular}{|p{2.8cm}|p{4.3cm}|p{4.2cm}|p{4.3cm}|}
\hline
\textbf{Bedarfskategorie} &
\textbf{Repräsentative User Needs} &
\textbf{Stakeholder Requirements} &
\textbf{Architekturwirkung} \\
\hline

Informationsqualität &
UN\_01\_001 Eignung von Bestandsquellen,
UN\_01\_003 Zuverlässigkeit,
UN\_01\_004 Vollständigkeit &
SR\_001 Input Data Suitability Assessment,
SR\_004 Source Reliability Classification &
Bewertung der Informationsgrundlage vor deren weiterführender
Engineering-Verarbeitung \\
\hline

Semantische Konsistenz &
UN\_01\_016 konsistentes technisches Vokabular,
UN\_01\_017 konsistente Interpretation semantisch äquivalenter Informationen,
UN\_01\_018 Erkennung ontologischer Konflikte &
SR\_016 Technical Vocabulary Alignment,
SR\_017 Semantic Equivalence Handling,
SR\_018 Ontology Conflict Detection &
Semantische Leitplanken für die Interpretation heterogener
Engineering-Informationen \\
\hline

Human Review und Intervention &
UN\_01\_012 Review generierter Interpretationen,
UN\_01\_013 Eingriff bei widersprüchlichen Interpretationen,
UN\_02\_003 Ergänzung fehlender oder widersprüchlicher Informationen &
SR\_012 Expert Review of Generated Interpretations,
SR\_013 Human Intervention for Conflicts,
SR\_026 Legacy Data Completion &
Explizite Einbindung menschlicher Fachentscheidungen in die
Informationsverarbeitung \\
\hline

Traceability und Governance &
UN\_01\_019 Rückverfolgbarkeit zur Ursprungsinformation,
UN\_04\_008 Zuordnung menschlicher Entscheidungen,
UN\_04\_011 Rekonstruktion der Verarbeitungshistorie &
SR\_019 Source Traceability,
SR\_035 Human Review Decision Association,
SR\_039 Processing History Traceability &
Erhalt von Herkunft, Verarbeitungskontext und Entscheidungsverlauf über
mehrere Verarbeitungsschritte hinweg \\
\hline

Projekt- und Source-Kontext &
UN\_01\_024 mehrere Engineering-Projekte,
UN\_01\_029 mehrere Quellen innerhalb eines Projekts,
UN\_04\_007 Trennung projektgebundener Informationen &
SR\_029 Multiple Project Support,
SR\_030 Project Information Isolation,
SR\_038 Multiple Source Support &
Explizite Trennung von Projekt- und Source-Kontexten sowie kontrollierte
Zusammenführung auf Projektebene \\
\hline

Modellsynthese und Modularität &
UN\_01\_006 Ableitung von Architekturinhalt,
UN\_01\_008 Wiederverwendbarkeit,
UN\_01\_009 nichtdestruktive Integration,
UN\_01\_010 modularer Modellzuwachs &
SR\_006 Architecture Information Derivation,
SR\_008 Model Reusability,
SR\_009 Modular Model Growth,
SR\_010 Non Destructive Submodel Integration &
Vorgaben für Struktur, Wiederverwendbarkeit und Integration erzeugter
Modellinhalte \\
\hline

Architecture as Code &
UN\_01\_022 maschinenverarbeitbare Architekturartefakte,
UN\_01\_023 SysML-v2-konforme Zielrepräsentation &
SR\_027 Architecture as Code Provision,
SR\_028 SysML v2 Target Representation &
Bereitstellung des resultierenden Architekturmodells als formale,
textuelle SysML-v2-Artefakte \\
\hline
\end{tabular}
\end{table}

Die in Tabelle~\ref{tab:architekturtreibende_requirements} dargestellten
Bedarfskategorien sind nicht unabhängig voneinander zu betrachten. Die
Verarbeitung mehrerer Quellen erfordert beispielsweise gleichzeitig eine
eindeutige Source-Zuordnung, die Trennung verschiedener Projektkontexte und die
Erhaltung der Traceability. Ebenso setzt die Erzeugung eines formalen
SysML-v2-Artefakts voraus, dass zwischen vorläufiger Interpretation und
fachlich autorisierter Engineering-Information unterschieden werden kann.

Besonders architekturprägend ist deshalb nicht ein einzelnes Requirement,
sondern das Zusammenspiel mehrerer Anforderungen. Aus diesem Zusammenspiel
entstehen Verantwortungsgrenzen, Informationszustände und Kontrollpunkte, die in
den nachfolgenden Abschnitten zur funktionalen und logischen Architektur
konkretisiert werden.

Die Nachvollziehbarkeit der Ableitung wird im Architekturmodell über die
verschiedenen Abstraktionsebenen erhalten. User Needs sind mit den daraus
abgeleiteten Stakeholder Requirements verbunden. Die Stakeholder Requirements
werden wiederum durch Use Cases adressiert und anschließend in die
Systemarchitektur überführt. Damit bleibt die Herkunft wesentlicher
Architekturentscheidungen über die Anforderungsebenen hinweg nachvollziehbar.
Die grundlegende Bedeutung einer solchen Traceability zwischen Ursprung,
Anforderung und nachgelagerten Engineering-Artefakten wurde bereits in
Abschnitt~\ref{sec:architekturmethodik} eingeordnet \cite{ISO29148}.

Die Use Cases ergänzen diese Anforderungssicht um stakeholderbezogene
Interaktionsszenarien. Sie beschreiben beispielsweise die Bewertung und
Aufbereitung von Bestandsinformationen, die Klärung von Mehrdeutigkeiten, die
Prüfung der Verwendbarkeit von Source-Beiträgen innerhalb eines Projekts sowie
die abschließende menschliche Bewertung eines generierten Modells. Sie sind
dabei noch nicht als interne Systemfunktionen oder als Festlegung von
Subsystemgrenzen zu verstehen, sondern beschreiben das erforderliche Verhalten
aus Sicht der beteiligten Stakeholder.

\subsection{Gesamtarchitektur und Engineering Information Flow}
\label{sec:engineering_information_flow}

Aus den in Abschnitt~\ref{sec:architektur_systemkontext} beschriebenen
Anforderungen wurde eine funktionale Verarbeitungskette abgeleitet, die den Weg
der Engineering-Information von einer registrierten Bestandsquelle bis zu einem
freigegebenen SysML-v2-Artefakt beschreibt. Der Informationsfluss bildet dabei
die zentrale Struktur der Informationsverarbeitungsarchitektur.

Die funktionale Betrachtung beschreibt zunächst, welche Verarbeitungsschritte,
Informationsübergänge und Kontrollentscheidungen erforderlich sind. Eine
Zuordnung dieser Funktionen zu logischen Komponenten oder Subsystemen erfolgt
erst in Abschnitt~\ref{sec:logische_architektur}. Damit bleibt die in
Abschnitt~\ref{sec:architekturmethodik} beschriebene Trennung zwischen
funktionaler und logischer Architekturebene erhalten.

Das Systemmodell unterscheidet hierzu einerseits einen übergeordneten
Engineering Transformation Lifecycle und andererseits detailliertere
Funktionszusammenhänge für einzelne Verarbeitungsbereiche. Der Lifecycle ist
nicht ausschließlich als lineare Ausführungssequenz zu verstehen. Neben der
Fortschreibung von Engineering-Information enthält er Review-, Gate- und
Rework-Pfade, durch die Verarbeitungsschritte erneut durchlaufen oder beendet
werden können.


\subsubsection{Funktionale Verarbeitungskette}
\label{sec:funktionale_verarbeitungskette}

Zu Beginn der Verarbeitung wird eine Engineering-Quelle einem eindeutig
bestimmten Projekt- und Source-Kontext zugeordnet. Dieser Kontext bildet die
Grundlage für die nachfolgende Verarbeitung und verhindert, dass
Engineering-Information ohne eindeutige Projekt- oder Source-Zugehörigkeit
weiterverwendet wird.

Für eine registrierte Quelle wird anschließend ein Verarbeitungskontext
erzeugt und die Source Information für die weitere Analyse vorbereitet.
Hierzu gehören die Bewertung ihrer grundsätzlichen Verarbeitbarkeit sowie die
Überführung in eine kontrollierte interne Repräsentation. Die ursprüngliche
Source bleibt dabei als Referenz erhalten.

Die im Entwicklungsverlauf getroffenen Architecture Decision Records (ADRs)
dokumentieren zentrale Architektur- und Realisierungsentscheidungen als
Engineering-Artefakte. Soweit sie für die Argumentation dieses Kapitels
relevant sind, werden sie im Folgenden als Design-Rationale-Evidence
herangezogen. Eine Übersicht sowie die für diese Arbeit relevanten ADR-Auszüge
sind in Anhang~\ref{app:architecture_decision_records} zusammengestellt.

Die Trennung zwischen Source Preparation und nachgelagerter semantischer
Interpretation wurde dabei als explizite Architekturgrenze festgelegt.
Dadurch bleibt die Aufbereitung der Ausgangsinformation von ihrer fachlichen
Interpretation getrennt und der Bezug zur ursprünglichen Engineering Source
über die weiteren Verarbeitungsschritte erhalten. Diese Entscheidung wurde
im Projekt unter anderem in ADR-009 und ADR-011 dokumentiert.

Auf dieser Grundlage erfolgt die source-lokale
Engineering-Informationsverarbeitung. Die vorbereitete Information wird in
analysierbare Informationseinheiten überführt und mit ihrer jeweiligen
Quellgrundlage verbunden. Aus diesen Einheiten wird
\emph{source-grounded Evidence} abgeleitet, die den Bezug zwischen späteren
Interpretationen und der zugrunde liegenden Source Information erhält.

Die explizite Trennung von Evidence Detection und Persona-basierter
Interpretation wurde im weiteren Architekturverlauf als eigene
Verarbeitungsgrenze konkretisiert. Source-grounded Evidence bildet dabei die
gemeinsame Nachweisbasis für nachgelagerte Interpretationen, ohne selbst bereits
eine fachliche Interpretation oder Freigabe darzustellen. Diese Ausgestaltung
ist in ADR-027 dokumentiert.

Die semantische Verarbeitung dient anschließend dazu, fachlich relevante
Engineering-Gegenstände zu identifizieren und unterschiedliche Benennungen oder
semantische Repräsentationen kontrolliert einzuordnen. Wo eine Interpretation
notwendig ist, können mehrere Persona-Perspektiven auf derselben Evidence-Basis
erzeugt und miteinander verglichen werden. Unterschiede zwischen diesen
Interpretationen werden nicht automatisch aufgelöst, sondern als zusätzliche
Bewertungsinformation erhalten.

Die theoretischen Grundlagen dieser Trennung zwischen Evidence,
probabilistischer Interpretation, Varianz und semantischer Governance wurden in
Abschnitt~\ref{sec:ki_gestaltungsprinzipien} beschrieben. In der Architektur
werden diese Prinzipien durch getrennte funktionale Verarbeitungsschritte
operationalisiert.

Das Ergebnis der source-lokalen Verarbeitung wird gemeinsam mit den zugehörigen
Evidence-, Bewertungs- und Traceability-Informationen für einen menschlichen
Engineering Review aufbereitet. Erst eine explizite fachliche Entscheidung
überführt ein geeignetes Zwischenergebnis in
\emph{Approved Engineering Information}. Die KI-gestützte Interpretation
besitzt damit selbst keine Engineering Approval Authority.

Bei Projekten mit mehreren Engineering-Quellen bleibt diese fachliche Freigabe
zunächst source-lokal. Erst nach Abschluss der jeweiligen Source-Verarbeitung
wird geprüft, welche freigegebenen Beiträge für den aktiven Projektkontext
zulässig sind. Dieser Übergang wird durch den \emph{Project Fit} kontrolliert.
Nur zugelassene Source-Beiträge dürfen in die projektweite Architekturableitung
einfließen.

Aus den zugelassenen Engineering-Informationen werden anschließend
Modellkandidaten abgeleitet. Die Modellbildung beinhaltet dabei neben der
fachlichen Ableitung auch die Entscheidung darüber, wie akzeptierte Inhalte in
die Zielmodellstruktur eingeordnet werden. Menschlich autorisierte
Modellkandidaten und Platzierungsentscheidungen werden danach in ein internes
Engineering-Modell überführt.

Erst auf Basis dieses kontrollierten internen Modells erfolgt die Erzeugung der
textuellen SysML-v2-Artefakte. Die eigentliche Generierung und die anschließende
Validierung der erzeugten Artefakte sind voneinander getrennte
Verarbeitungsschritte. Ein technisch valides Artefakt erhält dadurch noch keine
fachliche Modellautorität.

Die vollständige generierte Modellrevision wird deshalb abschließend gemeinsam
mit den verfügbaren Validierungsergebnissen und Traceability-Informationen einem
finalen Human Review unterzogen. Erst nach erfolgreicher technischer Prüfung und
expliziter fachlicher Autorisierung kann die betreffende Modellrevision für die
Veröffentlichung freigegeben werden.

% ABBILDUNG:
% CATIA View GV_EngineeringDataTransformationFlow
%
% Präsentationsorientierte End-to-End-Darstellung:
% Registered Engineering Source
% -> Process Information
% -> Human Review
% -> Project Fit
% -> Derive Architecture
% -> Validate Architecture
% -> Generate SysML
% -> Validate SysML
% -> Final Review
% -> Publication
%
% Die detaillierten internen Schritte von Process Information und
% Derive Architecture bleiben in dieser Übersicht bewusst zusammengefasst.
%
% \begin{figure}[htbp]
%     \centering
%     \includegraphics[width=\textwidth]{Abbildungen/...}
%     \caption{Funktionaler Engineering Information Flow des Turing Generators}
%     \label{fig:engineering_information_flow}
% \end{figure}

Abbildung~\ref{fig:engineering_information_flow} stellt diesen
End-to-End-Zusammenhang bewusst in verdichteter Form dar. Die detaillierten
Verarbeitungsschritte innerhalb der source-lokalen Informationsverarbeitung
sowie der Architekturableitung werden dabei zugunsten der Übersichtlichkeit
zusammengefasst.

\subsubsection{Einordnung des initialen Drei-Schichten-Modells}
\label{sec:three_layer_final_mapping}

Die in Abschnitt~\ref{sec:forschungsdesign} eingeführte frühe
Lösungshypothese unterschied die Informationsverarbeitung zunächst in einen
Data Layer, einen Process Layer und einen Knowledge Layer. Diese Einteilung
diente als konzeptioneller Ausgangspunkt der Architekturentwicklung und stellte
noch keine finale technische oder logische Zerlegung des Turing Generators dar.

In der ausgearbeiteten Architektur lassen sich die ursprünglichen drei
Betrachtungsebenen weiterhin erkennen, sie bilden jedoch keine voneinander
isolierten Architekturbereiche. Dem ursprünglichen Data Layer lassen sich
insbesondere die Verwaltung von Projekt- und Source-Kontexten, die Aufbereitung
der Ausgangsinformationen, die Bildung source-grounded Evidence sowie die
Aufrechterhaltung von Provenance zuordnen.

Der Process Layer findet sich in Funktionen zur Orchestrierung von
Verarbeitungsschritten, zur Steuerung von Processing States und Gates sowie in
der kontrollierten Weitergabe von Engineering-Information wieder. Dazu gehören
auch die Modellassemblierung, die SysML-v2-Generierung und die technische
Validierung.

Dem Knowledge Layer entsprechen insbesondere die semantische Governance, die
Interpretation durch unterschiedliche Persona-Perspektiven, die Bewertung von
Konsens und Varianz sowie Regeln für Architekturableitung und Model Placement.
Auch menschliche fachliche Entscheidungen beeinflussen den autorisierten
Bedeutungsgehalt der weiterverarbeiteten Engineering-Information.

Die finale Architektur zeigt jedoch, dass insbesondere Knowledge-,
Governance-, Traceability- und Authority-Funktionen querschnittlich über mehrere
Verarbeitungsphasen wirken. Eine starre Zuordnung der Architektur zu drei
technischen Schichten würde diese Zusammenhänge nicht angemessen abbilden.
Die weitere Darstellung folgt deshalb dem Engineering Information Flow und der
daraus abgeleiteten funktionalen und logischen Verantwortungsstruktur.

% =============================================================================
\subsection{Logische Architektur und Subsystemdekomposition}
\label{sec:logical_subsystems}
% =============================================================================

Aufbauend auf der funktionalen Architektur aus
Abschnitt~\ref{sec:functional_architecture} wird im nächsten Schritt
festgelegt, welche logischen Verantwortlichkeiten für die beschriebenen
Verarbeitungsfunktionen erforderlich sind und wie diese voneinander abgegrenzt
werden.

Die logische Zerlegung folgt dabei nicht unmittelbar der Reihenfolge des
Engineering Information Flow. Eine rein sequenzielle Zerlegung würde
Verarbeitungsphasen mit dauerhaft benötigten Verantwortlichkeiten vermischen.
Stattdessen werden zusammengehörige fachliche Verantwortungen in logisch
abgegrenzten Bereichen gebündelt. Funktionen, die an mehreren Stellen des
Informationsflusses benötigt werden, können dadurch einer gemeinsamen
logischen Verantwortung zugeordnet werden.

Die resultierende Architektur umfasst neun Subsysteme. Diese bilden die
übergeordneten Verantwortungsgrenzen der Zielarchitektur. Innerhalb der
Subsysteme werden die jeweiligen Anforderungen weiter in Subsystem Functions
und logische Komponenten konkretisiert. Die Subsystemstruktur ist damit eine
architektonische Verantwortungszerlegung und keine unmittelbare Abbildung
späterer Softwarepakete oder Implementierungsmodule.

\subsubsection{Verantwortungsbasierte Systemzerlegung}
\label{sec:verantwortungsbasierte_zerlegung}

% KERNAUSSAGEN:
%
% Die finale Architektur umfasst genau neun thesis-relevante Subsysteme:
%
% SBS_01 User Interaction and Guided Workflow
% SBS_02 Project and Source Context Management
% SBS_03 Processing Orchestration and State Control
% SBS_04 Engineering Information Processing
% SBS_05 Human Review and Engineering Authority
% SBS_06 Evidence Coverage and Traceability
% SBS_07 Architecture Derivation and Model Assembly
% SBS_08 SysML v2 Artifact Generation
% SBS_09 Generated Artifact Validation and Publication
%
% [MODEL/PROJECT]
% 2026-09-02 CATIA Modeling Complete SSOT.
%
% Finaler Modellumfang:
% 153 Subsystem Requirements
% 44 Subsystem Functions
% 31 Subsystem Logical Components
%
% Die Zahlen dienen als Umfangsnachweis.
% Sie sind NICHT die Argumentation des Abschnitts.
%
% ABBILDUNGSKANDIDAT 4.3:
% High-Level Logical Architecture mit neun Subsystemen.

Ausgangspunkt der logischen Zerlegung sind die auf Systemebene definierten
Funktionen. Diese beschreiben das erforderliche Verhalten unabhängig von einer
späteren technischen Realisierung. Die logische Architektur ordnet diese
Funktionen anschließend Verantwortungsbereichen zu, in denen zusammengehörige
Zustände, Entscheidungen und Informationsobjekte kontrolliert verarbeitet
werden.

Ein wesentliches Zerlegungskriterium ist dabei die Trennung unterschiedlicher
Authority- und Governance-Verantwortungen. Beispielsweise werden die
KI-gestützte Verarbeitung von Engineering-Information, menschliche
Reviewentscheidungen, Provenance und Traceability, Architekturableitung,
Artefakterzeugung und Artefaktvalidierung nicht in einer gemeinsamen
Verantwortung zusammengeführt.

Dadurch wird verhindert, dass eine Komponente gleichzeitig Informationen
interpretiert, deren fachliche Autorisierung festlegt und anschließend selbst
die technische Gültigkeit des daraus erzeugten Ergebnisses bestätigt. Die
logische Zerlegung unterstützt somit die in
Abschnitt~\ref{sec:informationszustaende_kontrollgrenzen} beschriebenen
Kontroll- und Authority-Grenzen.

Gleichzeitig bleiben projekt- und sourcebezogene Zustände von der eigentlichen
Informationsinterpretation getrennt. Orchestrierung und State Control bilden
ebenfalls eine eigenständige Verantwortung. Damit kann die Steuerung des
Verarbeitungsablaufs auf persistierten Zuständen und Entscheidungen basieren,
ohne selbst die fachliche Bedeutung der verarbeiteten Engineering-Information
zu bestimmen.

Die logische Architektur ist somit entlang stabiler Verantwortlichkeiten
strukturiert und nicht entlang einzelner Verarbeitungsschritte oder
Implementierungstechnologien.

Mehrere dieser Verantwortungsgrenzen wurden während der Architekturentwicklung
als eigenständige Architecture Decisions festgehalten. Dazu gehören
insbesondere die Trennung von source-lokaler Verarbeitung und Human Review,
die Einführung einer eigenständigen Model-Candidate-Ebene, die
representation-neutrale Internal-Engineering-Model-Grenze sowie die Trennung
von SysML-v2-Generierung, Artefaktvalidierung und finaler Human Authority.
Die zugehörigen Entscheidungen sind unter anderem in ADR-016, ADR-018,
ADR-019, ADR-021, ADR-022 und ADR-023 dokumentiert.

\subsubsection{Subsysteme und logische Verantwortlichkeiten}
\label{sec:subsysteme_logische_verantwortlichkeiten}

% TABELLEKANDIDAT 4.2:
% Verantwortlichkeiten der finalen Subsysteme
%
% Spalten:
% ID | Subsystem | zentrale Verantwortung | wichtigste Abgrenzung
%
% SBS_01 User Interaction and Guided Workflow
% - Human-facing Boundary
% - Guided Workflow
% - Status- und Evidence-Projektion
% - sammelt keine eigene Engineering Authority
%
% [PROJECT] ADR-017, ADR-024
%
% SBS_02 Project and Source Context Management
% - Projects und Sources registrieren / verwalten
% - Project Isolation
% - Source Selection / Source Roles / Context
% - entscheidet nicht ueber Engineering Meaning
%
% [PROJECT] ADR-005, ADR-012, ADR-032
%
% SBS_03 Processing Orchestration and State Control
% - Workflow Progression
% - Processing Run / Stage Control
% - kontrolliert nachfolgende Stufen
% - besitzt den finalen Publication Gate
% - erzeugt nicht selbst Engineering Meaning
%
% [PROJECT] finaler ownership boundary im 2026-09-02 SSOT
%
% SBS_04 Engineering Information Processing
% - Source Preparation
% - Evidence / Subjects / Semantic Processing
% - Project Fit Assessment
% - LLM-assisted interpretation where applicable
% - erzeugt Processing Evidence, nicht Engineering Approval
%
% [PROJECT] ADR-009, ADR-011, ADR-026, ADR-027, ADR-032
%
% SBS_05 Human Review and Engineering Authority
% - source-local Human Engineering Review
% - Approved Input / Approved Engineering Information
% - project-level Human Authority where required
% - Model Placement / Model Review Authority
% - Final Human Model Review
% - besitzt NICHT den technischen Publication Gate
%
% [PROJECT] ADR-016, ADR-023, ADR-029, ADR-032
%
% SBS_06 Evidence Coverage and Traceability
% - Evidence
% - Coverage
% - Provenance
% - Processing History
% - Traceability zwischen Source, Processing, Review, Model und Artifact
% - besitzt keine Engineering Decision Authority
%
% [PROJECT] ADR-012, ADR-027, ADR-032
%
% SBS_07 Architecture Derivation and Model Assembly
% - Project-level Model Candidate Set
% - Model Placement Integration
% - architecture derivation
% - deterministic Internal Engineering Model assembly
% - veraendert keine bereits autorisierte Engineering Semantics
%
% [PROJECT] ADR-018, ADR-019, ADR-028, ADR-029
%
% SBS_08 SysML v2 Artifact Generation
% - representation-neutral IEM in kontrollierte SysML-v2-Notation ueberfuehren
% - deterministic generation
% - exact profile-bound mappings
% - keine semantische Reinterpretation
%
% [PROJECT] ADR-021
%
% SBS_09 Generated Artifact Validation and Publication
% - generated-artifact validation
% - internal deterministic checks
% - external SYSIDE compatibility validation
% - blocking validation status
% - Name enthaelt "Publication", aber:
%   SBS_09 besitzt NICHT die Publication Authority.
%   Der finale Workflow-/Publication Gate liegt bei SBS_03,
%   Human Release Authority bei SBS_05.
%
% [PROJECT] ADR-022, ADR-023, 2026-09-02 SSOT
%
% KRITISCH:
% Diesen Punkt im finalen Text sauber erklaeren.

Aus dieser Zerlegung wurden neun Subsysteme abgeleitet. Sie decken gemeinsam
den vollständigen Engineering Information Flow ab, ohne dessen sequenzielle
Struktur unmittelbar als Systemstruktur zu reproduzieren.

\begin{table}[htbp]
\caption{Subsysteme der logischen Architektur und ihre zentralen Verantwortlichkeiten}
\label{tab:subsystem_verantwortlichkeiten}
\centering
\small
\begin{tabular}{|p{4.8cm}|p{9.2cm}|}
\hline
\textbf{Subsystem} &
\textbf{Zentrale Verantwortung} \\
\hline

SBS\_01 User Interaction and Guided Workflow &
Engineer-facing Interaktion, Projekt- und Source-Navigation,
Review-Interaktion sowie Darstellung von Processing Status und
Guided-Workflow-Informationen \\
\hline

SBS\_02 Project and Source Context Management &
Verwaltung persistenter Projektkontexte, eindeutiger Projekt- und
Source-Identitäten, Source-Registrierung sowie projektgebundener
Source-Auswahl \\
\hline

SBS\_03 Processing Orchestration and State Control &
Konfiguration und Koordination von Processing Runs, Steuerung von
Verarbeitungszuständen sowie Durchsetzung von Review-, Project-Fit- und
Publication-Gates \\
\hline

SBS\_04 Engineering Information Processing &
Source-lokale Aufbereitung, Evidence-Bildung, Engineering-Subject-Erkennung,
Persona-Interpretation und semantisch kontrollierte Verarbeitung von
Engineering-Information \\
\hline

SBS\_05 Human Review and Engineering Authority &
Durchführung und Persistierung menschlicher Engineering-Entscheidungen sowie
Etablierung der jeweiligen Human-Authority-Grenzen \\
\hline

SBS\_06 Evidence, Coverage and Traceability &
Bewertung von Evidence und Coverage sowie Erhaltung von Provenance,
Processing History und nachvollziehbarer Artefakt- und
Entscheidungszuordnung \\
\hline

SBS\_07 Architecture Derivation and Model Assembly &
Ableitung von Model Candidates aus autorisierter und für das Projekt
zugelassener Engineering-Information, Model Placement sowie deterministische
Assemblierung des Internal Engineering Model \\
\hline

SBS\_08 SysML v2 Artifact Generation &
Deterministische Überführung eines ausgewählten Internal Engineering Model
in ein strukturiertes und nachvollziehbares SysML-v2-Artifact Set \\
\hline

SBS\_09 Generated Artifact Validation and Publication &
Interne und externe Validierung des exakt erzeugten SysML-v2-Artifact Sets
sowie Ermittlung des technischen Publication-Blocking-Status \\
\hline
\end{tabular}
\end{table}

Die Subsystemgrenzen entsprechen damit unterschiedlichen Arten von
Verantwortung. SBS\_04 erzeugt und verarbeitet beispielsweise
Engineering-Interpretationen, besitzt jedoch keine menschliche
Freigabeautorität. Diese liegt in SBS\_05. SBS\_07 darf wiederum nur auf
Engineering-Information aufbauen, die bereits für die nachgelagerte
Architekturableitung autorisiert und für den betreffenden Projektkontext
zugelassen wurde.

Auch die Trennung von SBS\_08 und SBS\_09 ist bewusst gewählt. Die Erzeugung
einer SysML-v2-Repräsentation und deren Validierung sind unterschiedliche
Verantwortungen. Die Validierung bewertet ein bereits erzeugtes Artifact Set,
verändert dieses jedoch nicht. Ebenso begründet ein positives technisches
Validierungsergebnis allein keine Final Model Authority.

% ABBILDUNG:
% CATIA Cross-Subsystem Logical Architecture
%
% Darstellung der neun SBS als Verantwortungsblöcke und ihrer
% subsystemübergreifenden Interaktionen.
%
% \begin{figure}[htbp]
%     \centering
%     \includegraphics[width=\textwidth]{Abbildungen/...}
%     \caption{Subsystemdekomposition und subsystemübergreifende
%     logische Architektur des Turing Generators}
%     \label{fig:cross_sbs_logical_architecture}
% \end{figure}

Abbildung~\ref{fig:cross_sbs_logical_architecture} zeigt die neun Subsysteme
als logische Verantwortungsgrenzen sowie die wesentlichen Interaktionen zwischen
ihnen. Die Darstellung ergänzt damit den funktionalen End-to-End-Flow aus
Abschnitt~\ref{sec:functional_architecture} um die strukturelle Sicht auf die
Verantwortungsverteilung.

\subsubsection{Subsystemübergreifende Interaktionen und externe Architekturkontexte}
\label{sec:cross_sbs_external_contexts}

% KERNAUSSAGEN:
%
% 1. Die neun Subsysteme sind nicht als isolierte Pipeline-Bloecke zu verstehen.
%
% 2. Die finale Cross-SBS Architecture umfasst:
%    35 Logical Interfaces
%    56 typed Item Flows
%
% 3. Fuer die Thesis in drei Gruppen zusammenfassen:
%
% A. Control
%    - Orchestration steuert Processing, Review, Architecture, Generation
%      und Publication Progression.
%
% B. Engineering Information
%    - Source Context
%      -> Processing
%      -> Human Review
%      -> Approved Information
%      -> Architecture
%      -> SysML v2
%
% C. Evidence / Status
%    - wesentliche Stufen liefern Evidence und State
%      an SBS_06 bzw. SBS_01.
%
% 4. Vier Direction Gaps wurden vor Model Freeze korrigiert:
%    - Project / Source Context -> Evidence / Traceability
%    - Processing Evidence -> Human Review
%    - Coverage Assessment -> Architecture Derivation
%    - Processing State -> User Interaction / Status Presentation
%
% 5. Project Fit uebergibt sowohl:
%    ProjectFitAssessment
%    als auch ApprovedEngineeringInformation.
%
% [MODEL/PROJECT]
% 2026-09-02 CATIA closeout checkpoint.
%
% Keine Tabelle mit 35 Interfaces im Haupttext.

Die Subsysteme arbeiten nicht als voneinander isolierte Verarbeitungseinheiten.
Zwischen ihren logischen Grenzen werden definierte Informationen,
Entscheidungen, Zustände und Artefakte ausgetauscht. Die
Cross-Subsystem-Architektur beschreibt diese Interaktionen, ohne zusätzliche
fachliche Verantwortlichkeiten außerhalb der neun Subsysteme einzuführen.

Besonders deutlich wird dies am Engineering Information Flow. Projekt- und
Source-Kontext werden von SBS\_02 bereitgestellt, während SBS\_03 den
zugehörigen Processing Run und dessen Kontrollzustand koordiniert.
Source-lokale Verarbeitungsergebnisse aus SBS\_04 werden gemeinsam mit Evidence
und Traceability-Informationen aus SBS\_06 an die menschlichen
Review-Verantwortlichkeiten in SBS\_05 übergeben.

Nach erfolgreicher Autorisierung und Project-Fit-Prüfung können zugelassene
Engineering-Beiträge von SBS\_07 für die Architekturableitung verwendet
werden. Das resultierende Internal Engineering Model wird anschließend an
SBS\_08 übergeben und dort deterministisch in SysML-v2-Artefakte transformiert.
SBS\_09 bewertet schließlich das exakt erzeugte Artifact Set und stellt das
technische Validierungsergebnis für den abschließenden Review- und
Publication-Prozess bereit.

Neben diesen subsysteminternen Verantwortlichkeiten greift die Architektur auf
externe fachliche und technische Referenzkontexte zurück. Diese werden
ausdrücklich nicht als weitere Subsysteme modelliert, da sie keine
Engineering Authority innerhalb des Turing Generators besitzen.

Der \emph{LLM Inference Context} stellt die externe Inferenzfähigkeit für
ausgewählte KI-gestützte Verarbeitungsschritte bereit. Die erzeugten
Interpretationen bleiben jedoch Ergebnisse der Informationsverarbeitung und
erhalten durch den LLM-Aufruf keine fachliche Autorität.

Der Kontext \emph{Semantic Reference Knowledge} stellt kontrollierte
Terminologie- und Referenzinformationen für die semantische Verarbeitung
bereit. Dazu gehören insbesondere das Turing Core Vocabulary sowie verwendete
ontologische Referenzen. Diese Informationen unterstützen die semantische
Konsistenz, erzeugen oder autorisieren jedoch selbst keine
Engineering-Information.

\emph{Target Model and Framework References} stellen versionierte Regeln und
Zielstrukturen für Architekturableitung und Model Placement bereit. Sie
begrenzen damit den zulässigen Strukturraum der Modellableitung, ohne selbst
eine konkrete fachliche Platzierungsentscheidung zu treffen.

Der externe Kontext \emph{SysML v2 Modeling References} umfasst schließlich
kontrollierte Referenzen für Sprache, Syntax, Generierung und Validierung der
Zielrepräsentation. Dazu gehören die verwendete SysML-v2-Spezifikation,
validierte Syntaxreferenzen sowie ausgewählte Beispielmodelle. Diese
Referenzinformationen unterstützen die formale Repräsentation und technische
Prüfung, bleiben jedoch von der fachlichen Autorisierung des resultierenden
Modells getrennt.

Die Unterscheidung zwischen internen Verantwortlichkeiten und externen
Referenzkontexten verhindert damit insbesondere, dass die Nutzung eines
externen KI-Dienstes, einer Ontologie, eines Frameworks oder eines
Validierungswerkzeugs mit Engineering Authority gleichgesetzt wird.

\subsection{Evidence, Traceability und Human Authority}
\label{sec:evidence_authority}

Die in Abschnitt~\ref{sec:functional_architecture} beschriebene
Informationsverarbeitung erzeugt entlang des Engineering Information Flow
unterschiedliche Zwischen- und Endergebnisse. Für die Architektur ist dabei
nicht nur relevant, welche Information vorliegt, sondern auch, aus welchem
Kontext sie stammt, durch welche Verarbeitungsschritte sie entstanden ist und
welchen fachlichen Status sie besitzt.

Evidence, Traceability und Human Authority bilden deshalb einen
querschnittlichen Bestandteil der Architektur. Sie begleiten die
Informationsverarbeitung von der Engineering Source bis zum freigegebenen
Modellartefakt und verhindern, dass ein Verarbeitungsergebnis allein durch
seine technische Erzeugung oder Weitergabe fachliche Autorität erhält.

Die logische Verantwortung hierfür ist insbesondere zwischen
SBS\_05 \emph{Human Review and Engineering Authority} und
SBS\_06 \emph{Evidence Coverage and Traceability} aufgeteilt.
Während SBS\_05 menschliche Review- und Authority-Entscheidungen unterstützt
und verwaltet, erhält SBS\_06 die für diese Entscheidungen erforderliche
Evidence, Provenance und Verarbeitungshistorie.

\subsubsection{Informationszustände, Provenance und Traceability}
\label{sec:information_states_traceability}

Entlang des Engineering Information Flow werden unterschiedliche
Informationsobjekte verwendet, die jeweils eine bestimmte fachliche Rolle
innerhalb der Verarbeitung besitzen. Eine Source Projection oder
source-grounded Evidence hat beispielsweise einen anderen Status als
Approved Engineering Information oder ein Generated SysML Artifact Set.

Diese Unterscheidung ist erforderlich, da nicht jedes im System vorhandene
Informationsobjekt für nachgelagerte Engineering-Entscheidungen verwendet
werden darf. Insbesondere bleiben maschinell erzeugte Interpretationen,
Kandidaten und Bewertungsinformationen zunächst Evidence beziehungsweise
Vorschläge. Erst explizit autorisierte Informationsobjekte dürfen kontrollierte
Authority-Grenzen überschreiten.

Eine vereinfachte Folge zentraler Informationszustände lässt sich wie folgt
darstellen:

\begin{center}
\small
Engineering Source
$\rightarrow$ Source Projection
$\rightarrow$ Source-Grounded Evidence
$\rightarrow$ Engineering Subject
$\rightarrow$ Persona Interpretation
$\rightarrow$ Review Decision
$\rightarrow$ Approved Engineering Information
$\rightarrow$ Project Fit
$\rightarrow$ Model Candidate
$\rightarrow$ Approved Model Placement
$\rightarrow$ Internal Engineering Model
$\rightarrow$ Generated SysML Artifact Set
$\rightarrow$ SysML Validation Result
$\rightarrow$ Final Model Review Decision
$\rightarrow$ Publication Authorization
\end{center}

Diese Folge beschreibt keine zusätzliche Prozesssequenz neben
Abschnitt~\ref{sec:functional_architecture}. Sie verdeutlicht vielmehr die
unterschiedlichen Informations- und Authority-Zustände, die innerhalb des dort
beschriebenen Informationsflusses auftreten.

Für diese Zustände wird die Herkunft der Engineering-Information über die
Verarbeitung hinweg erhalten. Provenance umfasst dabei insbesondere den
zugehörigen Projekt- und Source-Kontext, den Processing Run und gegebenenfalls
den Processing Attempt sowie die Verbindung zwischen erzeugten Artefakten,
menschlichen Entscheidungen und nachgelagerten Modellzuständen.

Provenance dient damit nicht ausschließlich der nachträglichen Dokumentation.
Sie wird innerhalb der Architektur verwendet, um sicherzustellen, dass
nachgelagerte Verarbeitungsschritte auf den vorgesehenen und nachvollziehbaren
Informationsständen aufbauen. Dies betrifft insbesondere Human Review,
Project Fit, Model Placement, Modellassemblierung, Validierung und
Publication.

Die hierfür erforderliche Verarbeitungshistorie und Artefaktorganisation wurde
bereits früh als eigener persistenter Bestandteil der Architektur festgelegt.
ADR-012 dokumentiert hierzu die Trennung von Processing State,
Processing Evidence und versionierten Verarbeitungsartefakten als Grundlage
für die spätere End-to-End-Traceability.

SBS\_06 übernimmt hierzu sowohl die Erhaltung der Engineering Provenance als
auch die Verarbeitungshistorie und Artefakt-Governance. Die Traceability
umfasst dabei nicht nur die Verbindung eines Ergebnisses zu seiner
ursprünglichen Engineering Source, sondern auch die dazwischenliegenden
Informationsobjekte und Entscheidungen.

Eine resultierende Modellinformation kann dadurch prinzipiell über das
Internal Engineering Model, die zugrunde liegenden Model Candidates und die
Approved Engineering Information bis zu Reviewentscheidungen,
source-grounded Evidence und der ursprünglichen Engineering Source
zurückverfolgt werden.

Die in Abschnitt~\ref{sec:source_grounding} eingeführte Bedeutung von
Source Grounding wird damit auf die gesamte Verarbeitungskette erweitert.
Während Source Grounding zunächst den Nachweis zwischen Interpretation und
Ursprungsinformation sicherstellt, erhält die End-to-End-Traceability diesen
Bezug auch über nachgelagerte Engineering- und Modellierungsentscheidungen
hinweg.

% ABBILDUNGSKANDIDAT:
% Vereinfachte Authority- und Traceability-Kette
%
% Engineering Source
% -> Source-Grounded Evidence
% -> Human Review
% -> Approved Engineering Information
% -> Project Fit
% -> Model Candidate / Placement
% -> Internal Engineering Model
% -> Generated SysML Artifact Set
% -> Validation
% -> Final Model Review
%
% Unterhalb oder parallel dazu:
% Project / Source / Run / Attempt / Decision / Artifact Provenance
%
% \begin{figure}[htbp]
%     \centering
%     \includegraphics[width=\textwidth]{Abbildungen/...}
%     \caption{Informationszustände und Traceability entlang der
%     Authority-Grenzen des Turing Generators}
%     \label{fig:authority_traceability}
% \end{figure}

\subsubsection{Engineering Authority und Review-Grenzen}
\label{sec:authority_review_boundaries}

Die Architektur unterscheidet mehrere Formen menschlicher Authority, die an
unterschiedlichen Stellen des Engineering Information Flow benötigt werden.
Human Review ist deshalb nicht als einzelner Freigabeschritt am Ende der
Verarbeitung modelliert.

Eine erste Authority-Grenze besteht beim source-lokalen Human Engineering
Review. Dort wird bewertet, ob die aus einer Engineering Source abgeleitete und
interpretierte Information fachlich für eine weitere Engineering-Verwendung
akzeptiert werden kann. Erst durch eine explizite menschliche Entscheidung
kann daraus Approved Engineering Information entstehen.

Diese Engineering Authority bezieht sich zunächst auf den geprüften
Informationsinhalt und dessen Source-Kontext. Sie sagt noch nicht aus, dass die
Information für jedes Projekt oder jede konkrete Modellposition verwendet
werden darf.

Der Project Fit bildet deshalb eine hiervon getrennte Kontrollentscheidung.
Er bewertet, ob eine bereits reviewte source-lokale Information zum aktiven
Projektkontext zugelassen werden kann. Die source-lokale Engineering Authority
bleibt dabei erhalten. Der Project Fit ergänzt sie um die projektbezogene
Zulässigkeit.

Eine weitere Authority-Grenze entsteht bei der Architekturableitung.
Project-admitted Engineering Information kann zur Erzeugung von Model
Candidates verwendet werden, bestimmt jedoch nicht automatisch deren
Platzierung im Zielmodell. Die fachlich relevante Model Placement Decision
wird deshalb explizit durch Human Review autorisiert, bevor der betreffende
Inhalt in das Internal Engineering Model übernommen werden darf.

Die abschließende Authority-Grenze ist der Final Human Model Review. Dabei wird
nicht lediglich ein einzelnes Engineering-Element betrachtet, sondern die
konkrete generierte SysML-v2-Modellrevision gemeinsam mit ihrem
Validierungsergebnis, den Findings und der verfügbaren Traceability.

Ein positiver technischer Validation-Status stellt dabei keine menschliche
Engineering-Freigabe dar. Umgekehrt begründet eine menschliche
Engineering-Entscheidung allein noch keine technische Publication Eligibility.
Für die Veröffentlichung müssen die jeweils erforderlichen technischen und
menschlichen Voraussetzungen gemeinsam erfüllt sein.

Die Verantwortlichkeiten bleiben deshalb auch auf logischer Ebene getrennt.
SBS\_05 stellt die Human Review und Release Authority bereit. SBS\_09 bewertet
das erzeugte SysML-v2-Artefakt und liefert den zugehörigen Validation- und
Publication-Blocking-Status. SBS\_03 kontrolliert den finalen Publication Gate
auf Basis dieser Eingaben.

Der Final Human Model Review bezieht sich dabei auf die exakt überprüfte
Generated-SysML-v2-Artefaktrevision. Eine spätere inhaltliche Änderung führt
nicht dazu, dass die bestehende Freigabe automatisch auf den neuen Zustand
übertragen wird. Eine veränderte Revision muss erneut durch die hierfür
vorgesehenen Validierungs- und Review-Schritte geführt werden.

Damit entsteht eine gestufte Authority-Struktur. Automatisierte Verarbeitung
erzeugt Evidence, Interpretationen und Vorschläge. Menschliche Entscheidungen
autorisieren definierte Engineering- oder Modellzustände. Technische
Validierung prüft die resultierende Repräsentation. Erst das kontrollierte
Zusammenspiel dieser Verantwortlichkeiten ermöglicht die Freigabe eines
finalen Outputs.

\subsubsection{Identität, Immutability und Bindung von Entscheidungen}
\label{sec:identity_immutability}

Die Wirksamkeit der beschriebenen Authority-Grenzen setzt voraus, dass eine
Entscheidung eindeutig dem tatsächlich überprüften Informationszustand
zugeordnet werden kann. Eine Freigabe darf sich daher nicht lediglich auf einen
fachlichen Namen oder auf ein veränderbares Objekt beziehen.

Die Architektur trennt deshalb technische Identität und fachliche
Darstellungsbezeichnung. Projekt-, Source-, Processing- und Artefaktzustände
besitzen stabile technische Identitäten, über die ihre Beziehungen unabhängig
von späteren Änderungen ihrer Darstellung nachvollzogen werden können.

Gleichzeitig werden authority-relevante Artefakte und Entscheidungen als
immutable Zustände behandelt. Wird eine bereits geprüfte Engineering-Information
inhaltlich verändert, entsteht ein neuer Zustand, der die vorherige
Entscheidung nicht stillschweigend übernimmt. Historische Evidence und
Entscheidungen bleiben dadurch weiterhin ihrem ursprünglichen Gegenstand
zugeordnet.

Für die Bindung solcher Zustände werden zusätzlich Integritätsreferenzen
verwendet. Fingerprints ermöglichen es, Entscheidungen und nachgelagerte
Artefakte an einen konkreten Informationsstand zu koppeln. Sie dienen damit
der Identitäts- und Integritätsprüfung, sind jedoch kein Nachweis für die
fachliche Richtigkeit des Inhalts.

Besonders relevant wird diese Unterscheidung bei der Verarbeitung mehrerer
Engineering Sources. Eine lokale Artefaktidentität ist dort nicht zwingend
projektweit eindeutig. Für processing-run-bezogene Artefakte wird deshalb die
Kombination aus Processing Run, Artifact Type und Artifact Identity als
projektweit wirksamer Identitätskontext verwendet.

Die gleiche Logik gilt für menschliche Entscheidungen. Ein Human Review
Decision wird an den exakten Review-Gegenstand gebunden und kann nicht
implizit auf einen veränderten Nachfolgezustand übertragen werden. Dies gilt
für source-lokale Engineering Reviews ebenso wie für Model Placement und den
Final Human Model Review.

Im finalen Publication Gate werden diese Prinzipien zusammengeführt.
Publication Authorization hängt von dem konkreten Generated SysML Artifact
Set, dem dazugehörigen Validation Result und der korrespondierenden Human Final
Review Decision ab. Die Veröffentlichung bezieht sich damit auf einen
eindeutig bestimmten und überprüften Modellzustand und nicht lediglich auf
einen abstrakten Projekt- oder Modellnamen.

\subsection{Multi-Source-Verarbeitung und Modellableitung}
\label{sec:multisource_model_derivation}

Die Verarbeitung mehrerer Engineering Sources stellt eine besondere
Architekturanforderung dar, da unterschiedliche Quellen innerhalb desselben
Projekts nicht automatisch denselben fachlichen Geltungsbereich, dieselbe
Informationsqualität oder dieselbe Autorität besitzen.

Die Architektur führt deshalb mehrere Quellen nicht bereits während der
source-lokalen Verarbeitung zu einer gemeinsamen Engineering-Information
zusammen. Stattdessen wird zunächst für jede Source ein eigenständiger
Verarbeitungs- und Review-Kontext erhalten. Erst autorisierte und für den
jeweiligen Projektkontext zugelassene Source-Beiträge dürfen in die
projektweite Modellableitung eingehen.

Damit wird zwischen drei Fragestellungen unterschieden. Zunächst ist zu klären,
welche Engineering-Information aus einer einzelnen Source fachlich akzeptiert
werden kann. Anschließend wird ihre Zulässigkeit für das konkrete Projekt
bewertet. Erst danach wird untersucht, wie die zugelassenen Beiträge gemeinsam
in eine Zielmodellstruktur überführt werden können.


\subsubsection{Source-lokale Verarbeitung und Project Fit}
\label{sec:source_local_project_fit}

Die Verarbeitung einer Engineering Source erfolgt zunächst innerhalb ihres
eigenen Source- und Processing-Kontexts. Die in
Abschnitt~\ref{sec:functional_architecture} beschriebene Verarbeitung von
Source Information über Evidence, Interpretation und Human Engineering Review
wird damit zunächst unabhängig von weiteren im Projekt registrierten Sources
durchgeführt.

Das Ergebnis dieser Verarbeitung ist source-lokale
\emph{Approved Engineering Information}. Ihre fachliche Autorisierung bezieht
sich auf den konkret überprüften Informationsinhalt und dessen Source-Kontext.
Die bloße Registrierung einer Source innerhalb eines Engineering-Projekts
bedeutet noch nicht, dass ihre autorisierten Inhalte automatisch in die
projektweite Modellbildung eingehen dürfen.

Vor diesem Übergang wird deshalb ein \emph{Project Fit Assessment}
durchgeführt. Dabei werden die Approved Engineering Information, die
zugehörige Source Projection sowie Projekt-, Processing-Run- und
Processing-Attempt-Kontext gemeinsam betrachtet. Das Ergebnis beschreibt die
Zulässigkeit des Source-Beitrags für die weitere Verarbeitung innerhalb des
aktiven Projekts.

Nur Informationen, die diese Kontrollgrenze erfolgreich passieren, werden in
ein \emph{Project Admitted Engineering Information Set} aufgenommen. Project
Fit verändert dabei weder die ursprüngliche Source-Zuordnung noch die bereits
erteilte source-lokale Engineering Authority. Es ergänzt diese um die
projektbezogene Zulässigkeit.

Die Trennung ist insbesondere bei mehreren Sources relevant. Sie verhindert,
dass die Anwesenheit einer Source im Projekt bereits als fachliche
Gleichwertigkeit oder als automatische Freigabe für die gemeinsame
Modellbildung interpretiert wird. Ebenso wird keine allgemeine Rangordnung
zwischen Sources aus Source-Typ, Alter, Reihenfolge oder Anzahl abgeleitet.

Die einzelnen Source-Beiträge behalten auch nach Project Fit ihre jeweilige
Provenance und Authority-Zuordnung. Eine automatische semantische Verschmelzung
mehrerer source-lokaler Approved Engineering Information zu einer neuen
autorisierten Mischinformation ist an dieser Stelle nicht vorgesehen.

% ABBILDUNGSKANDIDAT:
% Multi-Source Project Fit / Admission
%
% Source A -> source-local processing -> Human Review -> Approved EI A \
%                                                               -> Project Fit
% Source B -> source-local processing -> Human Review -> Approved EI B /
%
% Project Fit
% -> Project Admitted Engineering Information Set
% -> Project-level Model Candidate Derivation
%
% Wichtig:
% Source-local Provenance und Authority bleiben sichtbar.
%
% \begin{figure}[htbp]
%     \centering
%     \includegraphics[width=\textwidth]{Abbildungen/...}
%     \caption{Source-lokale Verarbeitung und Project-Fit-basierte
%     Zulassung zur projektweiten Modellableitung}
%     \label{fig:multisource_project_fit}
% \end{figure}


\subsubsection{Projektweite Modellableitung und Human Model Placement}
\label{sec:project_model_derivation}

Nach der Project-Fit-Prüfung wechselt die Betrachtung von der source-lokalen
Informationsverarbeitung zur projektweiten Architekturableitung. Als Eingang
dienen ausschließlich die für das Projekt zugelassenen und bereits
Human-autorisierten Source-Beiträge.

Aus diesen Eingaben wird ein gemeinsames \emph{Model Candidate Set} für das
Projekt abgeleitet. Die Einführung einer eigenständigen Model-Candidate-Ebene trennt die
projektweite Architekturableitung von der bereits etablierten
Engineering Authority der Eingangsinformation. Diese Architekturentscheidung
ist in ADR-018 dokumentiert. Die zugrunde liegenden source-lokalen Authority- und
Provenance-Beziehungen bleiben dabei erhalten. Das Model Candidate Set stellt
somit eine abgeleitete Architekturrepräsentation dar und ersetzt seine
autorisierten Eingangsinformationen nicht.

Vor der eigentlichen Kandidatenbildung wird der verfügbare
Engineering-Informationsumfang hinsichtlich Coverage und Evidence Sufficiency
bewertet. Zusätzlich wird ein Model Derivation Context bestimmt, der die
anzuwendende Ableitungsstrategie und die Zielprojektion beschreibt.

Die Modellableitung kann deterministische und KI-gestützte Anteile kombinieren.
Wo eine eindeutige Zielprojektion aus den vorhandenen Regeln möglich ist, wird
diese bevorzugt deterministisch aufgelöst. Nicht eindeutig auflösbare
Projektionsentscheidungen bleiben als solche erhalten und dürfen nicht durch
implizite Annahmen verborgen werden.

Das erzeugte Model Candidate Set wird anschließend einem Human Candidate Review
zugeführt. Damit wird auch auf Projektebene zwischen maschinell beziehungsweise
regelbasiert abgeleiteter Architekturinformation und einem für die weitere
Modellbildung akzeptierten Kandidatenzustand unterschieden.

Für akzeptierte Kandidaten wird im nächsten Schritt eine Platzierung innerhalb
der gewählten Zielmodellstruktur vorgeschlagen. Model Placement beantwortet
dabei eine andere Frage als die vorherige Engineering-Interpretation:
Nicht mehr die fachliche Bedeutung der Source Information steht im Mittelpunkt,
sondern die Frage, an welcher Stelle und in welcher strukturellen Beziehung
diese bereits autorisierte Engineering-Information im Zielmodell repräsentiert
werden soll.

Placement-Vorschläge besitzen deshalb zunächst ebenfalls keine eigenständige
Model Authority. Erst der Human Model Placement Review erzeugt ein
\emph{Approved Model Placement Set}. Damit wird die konkrete strukturelle
Einordnung der Model Candidates explizit autorisiert. Die anschließende Model-Placement-Boundary folgt der in ADR-029 festgehaltenen
Entscheidung, dass eine strukturelle Platzierung vor der Modellassemblierung
explizit menschlich autorisiert werden muss.

Die Trennung zwischen Candidate Derivation und Model Placement verhindert,
dass die Ableitung fachlicher Architekturinformation unmittelbar mit einer
bestimmten Modellposition gleichgesetzt wird. Gleichzeitig bleibt die
Entscheidung über die Zielmodellstruktur nachvollziehbar von der
source-lokalen Engineering Authority getrennt.


\subsubsection{Internal Engineering Model, Generierung und Validierung}
\label{sec:iem_generation_validation}

Erst auf Basis eines vollständig autorisierten Model Candidate und Placement
State erfolgt die Assemblierung des \emph{Internal Engineering Model}.
Das Internal Engineering Model bildet eine representation-neutrale
Zwischenstufe zwischen fachlich autorisierter Modellinformation und ihrer
konkreten SysML-v2-Darstellung. Das Internal Engineering Model wurde entsprechend ADR-019 bewusst als
eigenständige, representation-neutrale Modellgrenze eingeführt. Es bildet
damit den autorisierten Engineering-Modellzustand unabhängig von seiner
späteren textuellen SysML-v2-Repräsentation ab.

Die Assemblierung erfolgt deterministisch auf Grundlage der autorisierten
Model Candidates, der freigegebenen Placement-Entscheidungen und der
ausgewählten Zielmodellstruktur. Während dieses Schrittes findet keine erneute
semantische Interpretation statt. Die zuvor etablierten Engineering- und
Model-Authority-Zustände werden damit in eine reproduzierbare interne
Modellrepräsentation überführt.

Vor der textuellen Repräsentation wird das Internal Engineering Model auf
architektonische Konsistenz und Integrationsfähigkeit geprüft. Dazu gehören
unter anderem die Einhaltung relevanter Architekturregeln,
Schnittstellenkonsistenz und die nichtdestruktive Integration in die
vorgesehene Modellstruktur.

Die anschließende SysML-v2-Generierung ist als deterministische
Repräsentationsfunktion ausgelegt. Ein ausgewählter Snapshot des Internal
Engineering Model wird unter Verwendung eines definierten
Generierungskontexts in ein maschinenverarbeitbares textuelles
SysML-v2-Artefakt überführt.

Die Generation darf dabei keine neue Engineering-Interpretation durchführen.
Sie überführt ausschließlich bereits repräsentierte und autorisierte Semantik
in die unterstützte Zielnotation. Nicht unterstützte Semantik darf daher nicht
stillschweigend durch eine andere Bedeutung ersetzt werden. Die in ADR-021 dokumentierte Generierungsarchitektur konkretisiert diese
Trennung, indem die SysML-v2-Erzeugung als deterministische Repräsentation
eines ausgewählten Internal-Engineering-Model-Zustands behandelt wird und
nicht als erneuter Interpretationsschritt.

Neben der textuellen Repräsentation wird auch die Struktur des resultierenden
Architecture-as-Code-Artefakts deterministisch festgelegt. Die erzeugten
Elemente und Beziehungen bleiben über maschinenauflösbare Traceability mit dem
Internal Engineering Model und den vorgelagerten autorisierten
Engineering-Informationen verbunden.

Die Validierung des Generated SysML Artifact Set ist von seiner Generierung
getrennt. Gegenstand der Prüfung ist der exakt erzeugte Artefaktstand. Die
Validierung darf diesen Zustand nicht verändern, sondern erzeugt
Validation Findings und einen zugehörigen Validation Status.

Die Architektur unterscheidet dabei interne Validierungsprüfungen von einer
externen Kompatibilitätsvalidierung. Interne Prüfungen adressieren unter
anderem für den Turing Generator definierte Integritäts-, Traceability-,
Struktur- und Relationship-Constraints. Die externe Validierung überprüft die
erzeugte Repräsentation zusätzlich gegen die verwendete SysML-v2-
Validierungsumgebung.

Ein erfolgreiches Validation Result stellt weiterhin keine Final Model
Authority dar. Wie in Abschnitt~\ref{sec:authority_review_boundaries}
beschrieben, erfolgt diese erst durch den Final Human Model Review auf Basis
des exakt generierten Artefaktstands und seines zugehörigen
Validierungsergebnisses. Erst danach kann der Publication Gate über die
Freigabe der betreffenden Revision entscheiden.

\subsection{Architekturgrenzen und Scope}
\label{sec:architecture_scope_boundaries}

Die in den vorangegangenen Abschnitten beschriebene Architektur bildet die
für die Forschungsfrage relevanten Requirements-, Functional- und
Logical-Aspekte des Turing Generators ab. Sie stellt jedoch weder eine
vollständige Produktarchitektur noch eine vollständige Beschreibung aller
für einen späteren produktiven Betrieb erforderlichen technischen Aspekte dar.

Die bewusste Abgrenzung des Architekturclaims ist erforderlich, um zwischen
dem konzipierten Zielbild, der prototypischen Realisierung und dem im Rahmen
der Arbeit tatsächlich untersuchten Funktionsumfang zu unterscheiden. Diese
Ebenen werden in den folgenden Kapiteln getrennt betrachtet.


\subsubsection{Physical und Deployment außerhalb des Thesis Scopes}
\label{sec:physical_out_of_scope}

Die in Abschnitt~\ref{sec:architekturmethodik} beschriebene
RFLP-Systematik dient als methodischer Bezugsrahmen für die
Architekturentwicklung. Im Rahmen dieser Arbeit werden jedoch bewusst nur die
Requirements-, Functional- und Logical-Ebene ausgearbeitet.

Eine Physical- oder Deployment-Architektur wird nicht detailliert modelliert.
Diese Entscheidung folgt aus dem Schwerpunkt der Forschungsfrage auf der
Informationsverarbeitungsarchitektur für die kontrollierte Überführung
heterogener Engineering-Information in SysML-v2-Strukturen.

Damit bleiben insbesondere konkrete Compute-Infrastrukturen,
Deployment-Topologien, Netzwerkarchitekturen, Runtime-Hosting und die
technische Verteilung einzelner Komponenten außerhalb des betrachteten
Architekturscopes.

Diese Abgrenzung bedeutet nicht, dass solche Aspekte für eine spätere
Produktivsetzung des Turing Generators ohne Bedeutung wären. Sie werden
lediglich nicht als Bestandteil des in dieser Arbeit untersuchten
Architekturproblems behandelt.

Die Logical Architecture beschreibt deshalb fachliche und technische
Verantwortungsgrenzen, ohne diese bereits auf konkrete Hosts, Prozesse,
Container, Cloud-Dienste oder andere physische Ausführungsstrukturen
abzubilden.


\subsubsection{Target Architecture und prototypisch realisierter Ausschnitt}
\label{sec:target_vs_mvp}

Das CATIA-Modell beschreibt die für die Thesis entwickelte
Engineering-Zielarchitektur des Turing Generators. Diese Zielarchitektur ist
nicht mit dem aktuell prototypisch implementierten Funktionsumfang
gleichzusetzen.

Einzelne Architekturverantwortlichkeiten und Requirements können daher als
Bestandteil des Zielbilds modelliert sein, obwohl sie innerhalb des
Proof of Concept nur teilweise oder noch nicht vollständig realisiert wurden.
Entscheidend ist, dass diese Abweichung zwischen Architektur und
Implementierungsstand explizit nachvollziehbar bleibt.

Für die weitere Arbeit werden deshalb drei Betrachtungsebenen unterschieden:
die modellierte Target Architecture, der tatsächlich implementierte
PoC-Ausschnitt und der im Rahmen der Verifikation und Validierung empirisch
untersuchte Ausschnitt.

Die Target Architecture beschreibt den konzeptionell vorgesehenen
Verantwortungs- und Informationsfluss. Kapitel~5 betrachtet dagegen, welche
Teile dieses Zielbilds im Turing Generator prototypisch umgesetzt wurden.
Kapitel~6 und Kapitel~7 grenzen anschließend ein, welche dieser realisierten
Funktionen tatsächlich Gegenstand der Verifikation und Validierung waren und
welche Ergebnisse dabei beobachtet wurden.

Diese Trennung verhindert, dass das Architekturmodell nachträglich auf den
jeweiligen Implementierungsstand reduziert wird. Gleichzeitig wird vermieden,
dass modellierte Zielanforderungen fälschlich als bereits implementiert oder
empirisch nachgewiesen interpretiert werden.


\subsubsection{Begrenzung des Architekturclaims}
\label{sec:architecture_claim_boundaries}

Das entwickelte Architekturkonzept beschreibt einen kontrollierten
Informationsverarbeitungspfad für den untersuchten Anwendungsfall und bildet
die Grundlage für die prototypische Realisierung. Aus der Architektur allein
lässt sich jedoch keine allgemeine Produktionsreife des Turing Generators
ableiten.

Insbesondere werden mit der vorliegenden Arbeit keine Aussagen über
produktive Skalierbarkeit, vollständiges Security Hardening,
Performance-Garantien oder die Eignung für beliebige
Deployment-Umgebungen getroffen. Ebenso wird keine vollständige Verarbeitung
beliebiger Medien-, CAD- oder Engineering-Datenformate beansprucht.

Die erzeugte SysML-v2-Repräsentation beschränkt sich auf den innerhalb der
Architektur und des Proof of Concept unterstützten Modellierungsumfang.
Daraus folgt weder eine vollständige Abdeckung sämtlicher SysML-v2-Konstrukte
noch eine allgemeine Übertragbarkeit auf beliebige Engineering-Domänen.

Auch die vorgesehenen Human-Review-, Evidence- und Governance-Mechanismen
stellen zunächst Architekturmechanismen dar. Ihre praktische Wirksamkeit und
ihre Grenzen müssen getrennt anhand der prototypischen Umsetzung und der
empirischen Ergebnisse bewertet werden. Diese Einordnung erfolgt in den
nachfolgenden Kapiteln und wird in Kapitel~8 im Hinblick auf die
Forschungsfragen und die Grenzen der Untersuchung diskutiert.

Zusammenfassend wurde aus der initialen Schichten- und Workflow-Hypothese eine
verantwortungsbasierte Requirements-, Functional- und Logical-Architecture
entwickelt. Die neun Subsysteme trennen die wesentlichen
Engineering-Verantwortlichkeiten voneinander. Engineering-Information
durchläuft dabei kontrollierte Informations- und Authority-Zustände, während
Provenance und Traceability ihre Herkunft und Verarbeitung über die
Architektur hinweg erhalten.

KI-gestützte Verarbeitung erzeugt innerhalb dieses Konzepts Evidence,
Interpretationen und Vorschläge, jedoch keine eigenständige Engineering
Authority. Fachliche Autorität entsteht an expliziten menschlichen
Review-Grenzen. Architekturableitung, Model Placement, interne
Modellassemblierung, SysML-v2-Generierung, technische Validierung und
Publication bleiben dabei als unterschiedliche Verantwortlichkeiten
voneinander getrennt.

Kapitel~5 betrachtet im nächsten Schritt, wie diese
Architekturverantwortlichkeiten im Turing Generator prototypisch realisiert
wurden.
